\documentclass[11pt]{article}
\usepackage[margin=1in]{geometry}
\usepackage{amsmath,amssymb}
\usepackage{booktabs}
\usepackage{hyperref}
\usepackage{siunitx}
\usepackage{xcolor}

\title{Independent Replication Report \\
  \large OSTI 1974586 --- Sewell, Bao, Jordan (2023) \\
  \emph{Variational quantum simulation of the critical Ising model
    with symmetry averaging}}
\author{Independent Replication (OSTI-100, rank 50)}
\date{2026-07-02}

\begin{document}
\maketitle

\begin{abstract}
We independently replicate the exactly-solvable reference physics that anchors
every quantitative claim in Sewell, Bao \& Jordan, \emph{Phys.\ Rev.\ A} \textbf{107},
042620 (2023) (OSTI 1974586, arXiv:2210.15053v2): the infinite-volume ground-state
energy density $-4/\pi$ of the 1D critical transverse-field Ising model (TFIM), the
free-fermion / Jordan--Wigner (matchgate) equivalence and the even-parity to
anti-periodic fermion boundary-condition mapping, the CFT ($c=1/2$) finite-size
scaling, the shallow-QAOA residual and the exact-preparation threshold at $p{=}L/2$
qubits, and the symmetry-averaging cancellation mechanism whose antiphase Kramers--Wannier
observable pair suppresses systematic error by $\sim$2 orders of magnitude.
The paper's headline \emph{algorithmic} contribution --- the full DMERA
matchgate-circuit optimization pipeline and its $\exp(-4.89\,D)$ scaling curves ---
was not rerun. \textbf{Verdict: PARTIAL.}
\end{abstract}

\section{Paper under test}
\begin{itemize}
  \item Sewell, Bao, Jordan, \emph{Phys.\ Rev.\ A} \textbf{107}, 042620 (2023).
  \item OSTI 1974586, DOI \texttt{10.1103/physreva.107.042620}, arXiv:2210.15053v2.
  \item OA PDF: \url{https://www.osti.gov/servlets/purl/1974586}
    (MD5 \texttt{131ff7c062bfb6993df7c222f7aaae49}).
\end{itemize}

\section{Paper summary}
The paper studies \emph{deep multi-scale entanglement renormalization} (DMERA)
quantum circuits as a variational ansatz for preparing critical ground states on
quantum computers, in the style of a VQE/VQA. The benchmark is the exactly-solvable
1D critical TFIM
\begin{equation}
  H_I = -\sum_j\!\bigl(X_j X_{j+1} + Z_j\bigr), \label{eq:HI}
\end{equation}
which is critical (CFT central charge $c=1/2$), integrable via Jordan--Wigner to a
free-fermion / Majorana model, and hence classically simulable to hundreds of qubits
using matchgate methods.

Key numerical claims:
\begin{enumerate}
  \item Infinite-volume ground-state energy density $= -4/\pi$ (Fig.~2 caption).
  \item Even-parity ground state (periodic spin BC) maps to \emph{anti-periodic}
    fermion boundary conditions ($\gamma_{2L+1} = -\gamma_1$).
  \item Energy-density error and state infidelity of the DMERA ansatz decay
    exponentially in depth $D$ with scaling coefficient $\approx -4.89$; DMERA
    outperforms a QAOA-style ansatz at equal depth.
  \item \emph{Symmetry averaging} over observables related by the (broken)
    Kramers--Wannier symmetry --- whose systematic errors are nearly out of phase
    --- cancels most of the error, reducing correlation-function error by
    $\sim 2$ orders of magnitude (relative error $<10^{-7}$ at $D=6$), and up to
    4 orders combined with translational averaging.
\end{enumerate}

\section{Claims table}
\begin{center}\small
\begin{tabular}{clllll}
\toprule
ID & Claim (abbreviated) & Type & Tested? & Result \\
\midrule
C1 & Infinite-volume $E/L = -4/\pi$ & scalar ref. & \checkmark & Reproduced $1.3\!\times\!10^{-13}$ \\
C2 & Free-fermion $\equiv$ spin; even-parity $\leftrightarrow$ ABC & equivalence & \checkmark & Reproduced ($<10^{-13}$) \\
C3 & Finite-$L$ $\to -4/\pi$ with $c{=}1/2$ FSS & convergence & \checkmark & Reproduced ($1/L^2$) \\
C4 & Shallow QAOA leaves large residual error & qualitative + threshold & \checkmark & Reproduced (5.8\% $p{=}1$; exact $p{=}4$) \\
C5 & KW-averaged antiphase obs. $\to \sim 2$ orders reduction & mechanism + magnitude & \checkmark & Mechanism reproduced (2.1 orders at $\sim 1^\circ$) \\
C6 & DMERA err $\propto \exp(-4.89\,D)$; beats QAOA/wavelet & full-circuit perf. & \ding{55} & \emph{Not rerun} (out of scope) \\
C7 & Up to 4 orders (translational $+$ KW) reduction & magnitude & \ding{55} & Consistent, not directly rerun \\
\bottomrule
\end{tabular}
\end{center}

\section{Method}
All computation in Python 3 / NumPy 2.x / SciPy (CPU; the exactly-solvable core is
tiny). No paper data reused beyond the analytic target $-4/\pi$ and Eq.~\eqref{eq:HI}.
\begin{itemize}
  \item \textbf{Download.} \texttt{curl} to \texttt{osti.gov} timed out on CherryRd;
    fetched via \texttt{ssh uicgpu} + \texttt{env.sh} proxy. Text via
    \texttt{pdftotext -layout} (born-digital PDF, no OCR needed).
  \item \textbf{C1/C2/C3 --- free-fermion spectrum.} For $J=h=1$ TFIM the
    single-particle dispersion is
    $\varepsilon(k) = 2\sqrt{(1-\cos k)^2 + \sin^2 k} = 4|\sin(k/2)|$;
    $E_0 = -\tfrac{1}{2}\sum_k \varepsilon(k)$ over even-sector (anti-periodic)
    momenta $k=(2m+1)\pi/L$. Infinite-volume density at $L=2\!\times\!10^6$.
    Cross-checked against brute-force dense diagonalization of the $L$-qubit spin
    Hamiltonian (PBC) via \texttt{numpy.linalg.eigvalsh} for $L=4,\ldots,12$.
  \item \textbf{C4 --- QAOA scan.} Standard TFIM-QAOA ansatz
    $\prod_{l=1}^{p} e^{-i\beta_l H_B} e^{-i\gamma_l H_C}\,|0\cdots 0\rangle$
    with $H_C = -\sum XX$, $H_B = -\sum Z$; energy minimized by Nelder--Mead with 10
    restarts at each $p$ on $L=8$ (PBC).
  \item \textbf{C5 --- symmetry-averaging mechanism.} Modeled two KW-related error
    signals $e_A, e_B$ that decay exponentially and are antiphase with residual phase
    mismatch $\phi$; verified the averaged-error suppression factor equals
    $|\sin(\phi/2)|$ (exact trig identity) and reproduces $\sim 2$ orders at
    $\phi \approx 1^\circ$.
  \item \textbf{Scoring.} Free Argo LLM-judge (\texttt{argo:gpt-5.2},
    \texttt{localhost:44497}); verdict PARTIAL.
\end{itemize}

\section{Results vs paper}

\subsection*{C1 --- infinite-volume energy density}
\begin{center}
\begin{tabular}{lccc}
\toprule
Quantity & This work & Paper ($-4/\pi$) & Abs.\ error \\
\midrule
$E_0/L$ ($L=2\!\times\!10^6$) & $-1.2732395447$ & $-1.2732395447$ & $1.3\!\times\!10^{-13}$ \\
\bottomrule
\end{tabular}
\end{center}

\subsection*{C2/C3 --- two methods, convergence to $-4/\pi$}
\begin{center}
\begin{tabular}{cccccc}
\toprule
$L$ & $E/L$ (free-fermion, ABC) & $E/L$ (dense spin, PBC) & $\|ff - \text{dense}\|$ & Err vs.\ $-4/\pi$ \\
\midrule
4    & $-1.30656296$ & $-1.30656296$ & $8.9\!\times\!10^{-16}$ & --- \\
8    & $-1.28145772$ & $-1.28145772$ & $1.8\!\times\!10^{-15}$ & $8.2\!\times\!10^{-3}$ \\
12   & $-1.27688293$ & $-1.27688293$ & $6.9\!\times\!10^{-14}$ & --- \\
64   & $-1.27336739$ & --- & --- & $1.3\!\times\!10^{-4}$ \\
256  & $-1.27324753$ & --- & --- & $8.0\!\times\!10^{-6}$ \\
1024 & $-1.27324004$ & --- & --- & $5.0\!\times\!10^{-7}$ \\
\bottomrule
\end{tabular}
\end{center}
Error scales as $\sim 1/L^2$ (halving $L$-spacing $\Rightarrow$ $\times 4$ reduction),
the expected $c=1/2$ CFT finite-size correction. Free-fermion and dense methods agree
to machine precision, confirming the Jordan--Wigner / matchgate description and the
even-parity $\leftrightarrow$ ABC mapping.

\subsection*{C4 --- QAOA-style shallow-circuit residual error ($L=8$, PBC)}
\begin{center}
\begin{tabular}{cccc}
\toprule
$p$ (depth) & $E_{\text{QAOA}}$ & $E_{\text{exact}}$ & Rel.\ energy error \\
\midrule
1 & $-9.656854$  & $-10.251662$ & $5.8\!\times\!10^{-2}$ \\
2 & $-9.952135$  & $-10.251662$ & $2.9\!\times\!10^{-2}$ \\
3 & $-10.054679$ & $-10.251662$ & $1.9\!\times\!10^{-2}$ \\
4 & $-10.251662$ & $-10.251662$ & $4.0\!\times\!10^{-13}$ \\
\bottomrule
\end{tabular}
\end{center}
Shallow circuits leave percent-level residual error (motivating DMERA); the exact
ground state is reached at $p=4$ for $2p=8$ spins, matching the paper's stated
exact-preparation threshold.

\subsection*{C5 --- symmetry-averaging suppression}
\begin{center}
\begin{tabular}{cccc}
\toprule
Residual phase $\phi$ & Single-obs max err & Averaged max err & Orders reduced \\
\midrule
$1^\circ$ & $5.6\!\times\!10^{-1}$ & $4.3\!\times\!10^{-3}$ & 2.12 \\
$2^\circ$ & $5.6\!\times\!10^{-1}$ & $8.5\!\times\!10^{-3}$ & 1.82 \\
$5^\circ$ & $5.6\!\times\!10^{-1}$ & $2.1\!\times\!10^{-2}$ & 1.42 \\
\bottomrule
\end{tabular}
\end{center}
At the near-antiphase regime the paper describes, KW averaging suppresses systematic
error by $\sim 2$ orders --- matching the paper's central quantitative claim.
Suppression factor $= |\sin(\phi/2)|$ (analytic).

\section{What was NOT reproduced}
\begin{itemize}
  \item \textbf{C6} --- the full DMERA matchgate circuit construction, its
    energy-minimization to $D \le 6$, and the $\exp(-4.89\,D)$ scaling curves
    (Figs.~1--4). These require the paper's appendix circuit parameters /
    optimization pipeline and are the paper's engineering contribution; out of
    scope for an efficient reference-physics replication.
  \item \textbf{C7} --- the full 4-orders combined (translational $+$ KW) reduction
    on actual DMERA observables. The mechanism is verified (C5); the end-to-end
    magnitude is not rerun.
\end{itemize}

\section{Genuine critique}
This is the deliberately harsh section --- what an adversarial reviewer would say to
both this replication effort \emph{and} the underlying paper.

\subsection*{Critique of this replication}
\begin{enumerate}
  \item \textbf{Only the ``easy half'' was reproduced.} Every claim we tick off (C1--C5)
    is a claim that a competent junior grad student could reproduce in a day: an
    exactly-solvable model, a Jordan--Wigner map, a QAOA sweep, and a two-line trig
    identity for the KW cancellation. The \emph{hard, novel} part of the paper is C6:
    building the DMERA circuit, encoding it as a matchgate network, and optimizing
    its Bogoliubov-scale parameters to hit the $\exp(-4.89\,D)$ curve. We did not touch
    that. Calling this a replication of ``the reference physics'' is honest, but the
    reference physics was never in doubt --- Onsager and Jordan--Wigner are 80+ years
    old. The interesting question (does DMERA really deliver the claimed scaling?)
    remains untested.
  \item \textbf{C5 is a mechanism check, not a numeric replication.} We showed that
    if two error signals are antiphase with residual mismatch $\phi$, then
    $|\sin(\phi/2)|$ gives $\sim 2$ orders at $\phi \approx 1^\circ$. That is a
    tautology of trigonometry, not evidence that the paper's actual KW-observable
    pair really lives near $\phi \approx 1^\circ$ on real DMERA output. The paper's
    $10^{-7}$ figure at $D=6$ could still be wrong even if our sanity check passes.
  \item \textbf{L=8 is small.} The QAOA scan was performed at $L=8$ where the ansatz
    at $p=L/2=4$ is trivially exact by parameter counting. That is the boundary
    condition where any QAOA-style ansatz becomes exact --- reproducing that
    threshold is unsurprising and does not test the paper's claim that DMERA still
    outperforms QAOA at $L \gg 2p$.
  \item \textbf{No paper-code diff.} We did not attempt to obtain and re-execute the
    authors' matchgate optimization pipeline (if released). A ``true'' partial
    replication would at minimum audit their optimization loss functions, ansatz
    layout, and reported minima; we did not.
  \item \textbf{Single-judge scoring.} The LLM judge is a single Argo model
    (\texttt{argo:gpt-5.2}) without cross-model disagreement gating. A one-vote
    concurrence with our own PARTIAL verdict is weak independent evidence.
  \item \textbf{Downloading via a proxy is a minor provenance risk.} The PDF was
    fetched via \texttt{ssh uicgpu} because CherryRd timed out on \texttt{osti.gov};
    the MD5 was recorded but not cross-verified against a second source (arXiv PDF
    MD5). Not a serious concern for a well-known preprint, but worth logging.
\end{enumerate}

\subsection*{Critique of the paper}
\begin{enumerate}
  \item \textbf{Symmetry averaging generalization is asserted, not shown.} The paper
    argues KW-averaging works because the KW-related observables have systematic
    errors that are nearly out of phase. That is elegant when it holds, but the
    mechanism is model-specific. Whether analogous ``antiphase'' pairings exist for
    non-Abelian symmetries, higher-dimensional CFTs, or non-integrable critical
    models is a live open question the paper does not answer.
  \item \textbf{Error-model assumptions are optimistic.} The reported error curves
    are computed on ideal classical simulators. On real NISQ hardware with coherent
    crosstalk and non-Markovian noise, the two KW-related observables may no longer
    have anti-correlated errors --- suppression could degrade or invert. The paper
    does not include a hardware-noise stress test.
  \item \textbf{Statistical vs.\ systematic error trade-off.} Averaging over two
    observables cancels systematic error at the cost of $\sqrt{2}$ statistical
    variance per shot. At $10^{-7}$ correlation-function precision, the shot budget
    required to actually measure the KW-averaged result is enormous and not
    discussed.
  \item \textbf{Comparison to mid-circuit measurement / post-selection.} An
    alternative to symmetry-averaging is to enforce the symmetry via mid-circuit
    measurement and post-select. The paper does not benchmark against that
    alternative.
  \item \textbf{$-4.89$ scaling: no error bar.} The reported scaling coefficient
    $\approx -4.89$ has no reported uncertainty. Is this a fit over 3 depth points or
    30? What is the confidence interval? Without that, ``exponential scaling with
    slope $-4.89$'' is a curve-fit assertion, not a verified law.
  \item \textbf{Definition of ``beats wavelet ansatz'' unclear.} The comparison to
    wavelet-style circuits is made at ``equal depth,'' but wavelet ansätze differ in
    gate count per layer, connectivity, and parameter count --- ``depth'' alone is an
    ambiguous axis.
\end{enumerate}

\section{Assessment and verdict}
The exactly-solvable reference physics that anchors every quantitative claim in the
paper is independently confirmed with high confidence. The paper's headline
\emph{algorithmic} contribution (DMERA circuit optimization + full performance
curves) was not rerun. Per the wave rubric this is a clear \textbf{PARTIAL}: core
reference physics and mechanisms replicated; end-to-end DMERA performance unverified.
The independent LLM-judge concurred.

\medskip
\noindent \textbf{Verdict: PARTIAL.}

\end{document}
